% v2.7.0 figure UID: FIG-P736-01
% Four aligned teaching routes plus two short, model-specific reuse examples.
\tikzset{slfig-FIG-P736-01/.style={font=\fontsize{9.6pt}{11.6pt}\selectfont,every path/.append style={line cap=round,line join=round}}}
\pgfkeysifdefined{/tikz/alt}{}{\tikzset{alt/.code={}}}
\begin{figure}[H]
\centering
\begin{tikzpicture}[slfig-FIG-P736-01,
  every node/.style={font=\fontsize{9.6pt}{11.6pt}\selectfont,align=center,outer sep=0pt},
  family/.style={draw=SLMidBlue,fill=SLSoftBlue,line width=.78pt,
    rounded corners=2.5pt,minimum width=27mm,minimum height=11mm,
    inner xsep=1.8mm,inner ysep=1.2mm},
  engine/.style={draw=SLRuleGray,fill=SLSoftGray,line width=.74pt,
    rectangle,minimum width=27mm,minimum height=11mm,
    inner xsep=1.8mm,inner ysep=1.2mm},
  rowtitle/.style={font=\bfseries\fontsize{10.2pt}{12.2pt}\selectfont,
    text=SLDeepBlue,inner sep=0pt},
  legendbox/.style={draw=SLRuleGray,fill=white,line width=.72pt,
    rounded corners=2pt,text width=135mm,
    font=\fontsize{9.6pt}{11.6pt}\selectfont,
    inner xsep=2mm,inner ysep=1.5mm},
  routeexample/.style={font=\fontsize{9.6pt}{11.6pt}\selectfont,
    fill=white,text=SLTextGray,inner sep=1pt},
  mainroute/.style={->,draw=SLDeepBlue,line width=.96pt},
  optional/.style={->,draw=SLRuleGray,line width=.78pt,densely dashed},
  >={Stealth[length=2.2mm,width=1.55mm]},
  alt={对象—关系—结论：上层依次是降维、话题分析、聚类和图分析四个方法族，下层与之对齐的是SVD、后验推断、EM和幂法四个计算引擎。四条竖直实线只表示本讲义选取的代表性主路线；话题分析到SVD和EM的两条短虚线分别以LSA和PLSA为例，表示特定模型或求解方案中的条件性复用。线型只区分关系类别，不编码强弱；实线与虚线都不是方法族和引擎之间的必要或充分关系，共享某个引擎也不使方法等价。方法族用圆角蓝底框，引擎用直角灰底框，主路线用实线、条件性复用用虚线，并由文字图例重复编码。}]
\node[rowtitle] (familytitle) at (0,3.15) {方法族};
\matrix (top) [matrix of nodes,column sep=8mm,nodes={family}] at (0,1.65)
  {\node (dim) {降维}; & \node (topic) {话题分析}; &
   \node (cluster) {聚类}; & \node (graph) {图分析}; \\};
\matrix (bottom) [matrix of nodes,column sep=8mm,nodes={engine}] at (0,-1.15)
  {\node (svd) {SVD}; & \node (posterior) {后验推断}; &
   \node (em) {EM}; & \node (power) {幂法}; \\};
\node[rowtitle,text=SLTextGray] (enginetitle) at (0,-2.45) {共享计算引擎};
\node[legendbox] (legend) at (0,-3.95)
  {\textbf{实线：}本讲义中的代表性主路线\qquad
   \textbf{虚线：}特定模型/求解方案中的条件性复用\\
   线型只分关系类型，不编码强弱；两类边均非必要或充分关系；\\
   共享引擎不意味着方法等价。};
\draw[mainroute] (dim)--(svd);
\draw[mainroute] (topic)--(posterior);
\draw[mainroute] (cluster)--(em);
\draw[mainroute] (graph)--(power);
\draw[optional] (topic.south west)--(svd.north east)
  node[midway,above=1mm,routeexample] {如LSA};
\draw[optional] (topic.south east)--(em.north west)
  node[midway,above=1mm,routeexample] {如PLSA};
\end{tikzpicture}
\captionsetup{width=.94\linewidth}
\caption{方法族—计算引擎教学映射：实线为本讲义代表性主路线，虚线为特定模型/方案的条件性复用；这些边均非必要或充分关系，共享引擎也不使方法等价。}
\label{fig:V5-C08-method-map}
\end{figure}
